%! The Nullspace of A: Solving Ax = 0 — Unit 3, Lesson 3.2 — Guided Notes (Answer Key)
\documentclass[10pt]{article}
\usepackage{linalg-article}
\usepackage{linalg-key}   % pulls in -boxes; do NOT also load -boxes
\newcommand{\TallMath}[1]{$\displaystyle #1\rule[-1.4em]{0pt}{3.2em}$}
\newcommand{\vv}{\mathbf{v}}
\newcommand{\ww}{\mathbf{w}}
\newcommand{\xx}{\mathbf{x}}
\newcommand{\yy}{\mathbf{y}}
\newcommand{\bb}{\mathbf{b}}
\newcommand{\svec}{\mathbf{s}}
\newcommand{\zero}{\mathbf{0}}

\begin{document}

\pageheader{Unit 3, Lesson 3.2}{Guided Notes --- Answer Key}
\namedateperiod

\vspace{0.12in}

\begin{objectivebox}
\textbf{By the end of this lesson, I will be able to\dots}
\begin{itemize}\setlength{\itemsep}{2pt}
  \item describe the \emph{nullspace} $N(A)$ --- all solutions of $A\xx=\zero$ --- and say why it is a subspace;
  \item solve $A\xx=\zero$ by \emph{elimination}, telling \emph{pivot columns} from \emph{free columns};
  \item build a \emph{special solution} for each free variable and describe $N(A)$ as a point, line, or plane.
\end{itemize}
\end{objectivebox}

\vspace{0.10in}

\begin{vocabbox}
\small
Fill in each term as we build it together.
\vspace{2pt}
\vocabans{Nullspace $N(A)$}{all solutions $\xx$ of $A\xx=\zero$; a subspace of $\mathbb{R}^n$ ($n$ = number of columns)}
\vocabans{Pivot column / pivot variable}{a column holding a pivot after elimination; its variable is solved for (not free)}
\vocabans{Free column / free variable}{a column with no pivot; its variable is free to choose any value}
\vocabans{Special solution}{the nullspace vector from setting one free variable to $1$, the rest to $0$, then solving back}
\vocabans{Free variables $=n-r$}{columns minus pivots; this is the number of special solutions = dimension of $N(A)$}
\end{vocabbox}

\begin{hookbox}
A robot arm has two levers; the setting $\xx=\begin{bsmallmatrix} x_1\\x_2 \end{bsmallmatrix}$ moves the
tip to $A\xx$. For $A=\begin{bsmallmatrix} 1 & 2\\ 2 & 4 \end{bsmallmatrix}$ the two lever directions
$\begin{bsmallmatrix} 1\\2 \end{bsmallmatrix}$ and $\begin{bsmallmatrix} 2\\4 \end{bsmallmatrix}$ lie on
the \emph{same} line.
\begin{itemize}\setlength{\itemsep}{1pt}
  \item Find one nonzero setting $\xx$ that leaves the tip unmoved ($A\xx=\zero$). \ansline{$\xx=\begin{bsmallmatrix} -2\\1 \end{bsmallmatrix}$: $-2\begin{bsmallmatrix} 1\\2 \end{bsmallmatrix}+1\begin{bsmallmatrix} 2\\4 \end{bsmallmatrix}=\zero$ (push lever~2 forward $1$, pull lever~1 back $2$).}
  \item Is there more than one such setting? Describe them all. \ansline{Yes --- every multiple $t\begin{bsmallmatrix} -2\\1 \end{bsmallmatrix}$ works. They form a line through the origin: the nullspace.}
\end{itemize}
\end{hookbox}

\begin{notesbox}{1. The Nullspace --- and Why It Is a Subspace}
The \textbf{nullspace} of $A$, written $N(A)$, is the set of \emph{all} solutions $\xx$ of
$A\,\xx=\ans{$\zero$}$. It always contains the zero vector, because $A\zero=\zero$. Run the Lesson~3.1
subspace requirement on it: if $A\xx_1=\zero$ and $A\xx_2=\zero$, then
\[
  A(\xx_1+\xx_2)=A\xx_1+A\xx_2=\zero, \qquad A(c\,\xx_1)=c\,(A\xx_1)={\color{keyred}\mathbf{\zero}}.
\]
Both closure rules hold, so $N(A)$ is a \textbf{subspace} of $\mathbb{R}^{\,{\color{keyred}\mathbf{n}}}$ (here $n$ =
number of \emph{columns} of $A$, because $\xx$ has one entry per column). \emph{Contrast:} the column
space $C(A)$ lived in $\mathbb{R}^m$ (outputs); the nullspace lives in $\mathbb{R}^n$ (inputs).
\end{notesbox}

\begin{notesbox}{2. Solve $A\xx=\zero$ by Elimination: Pivot vs.\ Free Columns}
The right-hand side is all zeros, and stays $\zero$ under \emph{every} row operation --- so to solve
$A\xx=\zero$ we just row-reduce $A$. Take $A=\begin{bmatrix} 1 & 1 & 2 \\ 1 & 2 & 3 \end{bmatrix}$:
\[
  \begin{bmatrix} 1 & 1 & 2 \\ 1 & 2 & 3 \end{bmatrix}
  \xrightarrow{\;R_2-R_1\;}
  \begin{bmatrix} 1 & 1 & 2 \\ 0 & 1 & 1 \end{bmatrix}
  \xrightarrow{\;R_1-R_2\;}
  \begin{bmatrix} 1 & 0 & {\color{keyred}\mathbf{1}} \\ 0 & 1 & {\color{keyred}\mathbf{1}} \end{bmatrix}.
\]
The pivots sit in columns $1$ and $2$ --- these are the \textbf{pivot columns}, and $x_1,x_2$ are
\textbf{pivot variables}. Column $3$ has \emph{no} pivot: it is a \textbf{free column}, and $x_3$ is a
\textbf{free variable} --- we are free to \ans{choose} its value. Reading the reduced rows back:
\[
  x_1 + x_3 = 0, \qquad x_2 + x_3 = 0.
\]
\end{notesbox}

\begin{notesbox}{3. Special Solutions --- and the Whole Nullspace}
To pin down a nullspace vector, give the free variable the value \ans{$1$} and solve back for the
pivot variables. With $x_3=1$: $\;x_1=-x_3=\ans{$-1$}$ and $x_2=-x_3=\ans{$-1$}$, so the
\textbf{special solution} is
\[
  \svec=\begin{bmatrix} x_1 \\ x_2 \\ x_3 \end{bmatrix}
  =\begin{bmatrix} \rule{0pt}{1.0em}{\color{keyred}\mathbf{-1}} \\ {\color{keyred}\mathbf{-1}} \\ {\color{keyred}\mathbf{1}} \end{bmatrix}.
\]
Every solution of $A\xx=\zero$ is a multiple of $\svec$, so $N(A)=\{\,t\,\svec : t \text{ any number}\,\}$.
\textbf{The rule:} there is one special solution \emph{per free variable}, and the nullspace is all
\ans{combinations} of the special solutions. The number of free variables is
\[
  n-r \;=\; (\text{columns}) - (\text{pivots}) \;=\; 3 - 2 \;=\; {\color{keyred}\mathbf{1}}.
\]
\end{notesbox}

\begin{notesbox}{4. Point, Line, or Plane}
The free count $n-r$ is the nullspace's \emph{dimension}, so it tells you the shape:
\begin{itemize}[leftmargin=1.5em, itemsep=2pt]
  \item $n-r=0$ (no free variables): only $\xx=\zero$, so $N(A)=\{\zero\}$ --- a \textbf{point}.
  \item $n-r=1$: one special solution, so $N(A)$ is a \textbf{line} through the origin (our example above).
  \item $n-r=2$: two special solutions, so $N(A)$ is a \textbf{plane} through the origin.
\end{itemize}
A plane example: $A=\begin{bmatrix} 1 & 2 & 3 \end{bmatrix}$ gives one equation $x_1+2x_2+3x_3=0$ with
pivot column $1$ and free columns $2,3$. Setting $(x_2,x_3)=(1,0)$ then $(0,1)$:
\[
  \svec_1=\begin{bmatrix} \rule{0pt}{1.0em}{\color{keyred}\mathbf{-2}} \\ 1 \\ 0 \end{bmatrix}, \qquad
  \svec_2=\begin{bmatrix} \rule{0pt}{1.0em}{\color{keyred}\mathbf{-3}} \\ 0 \\ 1 \end{bmatrix},
\]
and $N(A)$ is the plane of all combinations $t_1\svec_1+t_2\svec_2$ in $\mathbb{R}^3$ ($n-r=3-1=2$).

\vspace{3pt}
\textbf{Next (3.3):} attach a right side --- the \emph{complete} solution of $A\xx=\bb$ is one
particular solution \emph{plus} the whole nullspace.
\end{notesbox}

\begin{practicebox}
Show your reasoning.
\begin{enumerate}[leftmargin=1.6em, itemsep=4pt]
  \item Is $\xx=\begin{bsmallmatrix} 2\\-1 \end{bsmallmatrix}$ in the nullspace of $A=\begin{bsmallmatrix} 1 & 2 \\ 3 & 6 \end{bsmallmatrix}$? Check $A\xx$. \ansline{Yes --- $A\xx=\begin{bsmallmatrix} 1\cdot2+2\cdot(-1)\\3\cdot2+6\cdot(-1) \end{bsmallmatrix}=\begin{bsmallmatrix} 0\\0 \end{bsmallmatrix}=\zero$.}
  \item Find the nullspace of $A=\begin{bsmallmatrix} 1 & 3 \\ 2 & 6 \end{bsmallmatrix}$: give the special solution and say whether $N(A)$ is a point, line, or plane. \ansline{$R_2-2R_1$ gives $\begin{bsmallmatrix} 1&3\\0&0 \end{bsmallmatrix}$, so $x_1+3x_2=0$ with $x_2$ free. Special solution $\svec=\begin{bsmallmatrix} -3\\1 \end{bsmallmatrix}$; $N(A)$ = all multiples of $\svec$ = a \emph{line} through the origin ($n-r=2-1=1$).}
  \item A matrix has $5$ columns and rank $r=3$. How many free variables does $A\xx=\zero$ have? What is the dimension of $N(A)$? \ansline{Free variables $=n-r=5-3=2$; so $N(A)$ has $2$ special solutions and dimension $2$ (a plane through the origin, inside $\mathbb{R}^5$).}
\end{enumerate}
\end{practicebox}

\begin{teachernote}
The two must-dos are \S2--\S3 (find pivots/free columns, then build the special solution by setting the
free variable to $1$ and solving back) and the home rule in \S1 ($N(A)\subseteq\mathbb{R}^n$, not
$\mathbb{R}^m$ --- it is a set of \emph{inputs}). The most common slip is stopping at one special
solution: stress ``one per free variable, then take all combinations.'' If time is short, \S4's plane
example can be assigned rather than worked live.
\end{teachernote}

\end{document}
